\documentclass[11pt,a4paper]{article}
\usepackage[utf8]{inputenc}
\usepackage{geometry}
\geometry{margin=1in}
\usepackage{lineno}
\usepackage{xcolor}
\usepackage{titlesec}
\usepackage{array}

% Command for redacting personal information with a black box
\newcommand{\redacted}{\rule{2.5cm}{2ex}}

\title{\textbf{Social Science Fair Interview 03 \\ Transcript: Value and Pressure of Football Defenders (Right-Back Perspective II)}}
\author{\textbf{Research Team Members:} Luke Sun, Kiara Wei, Winston Huang}
\date{May 28, 2026}

\begin{document}

\maketitle

\section*{Transcript Metadata}
\begin{tabular}{ll}
    \textbf{Project Theme:} & Discussion on the Value and Pressure of Football Defenders \\
    \textbf{Date of Interview:} & Thursday, May 28, 2026 \\
    \textbf{Time of Interview:} & 16:01 -- 16:17 (GMT+08) \\
    \textbf{Interviewer Identifier:} & Speaker 1 Luke Sun\\
    \textbf{Interviewee Identifier:} & Speaker 2 (\redacted) \\
    \textbf{Interviewee Subject Profile:} & School Team Athlete (Right-Back / Full-Back)
\end{tabular}

\vspace{2em}
\hrule
\vspace{1.5em}

\section*{Official Verbatim Transcript}

\begin{linenumbers}

\noindent \textbf{Speaker 1 [00:04]:} Okay, our interview is starting now. \redacted, you are on the school team, right? Football.

\vspace{0.5em}
\noindent \textbf{Speaker 2 [00:17]:} Defender. Defender. Ah wait, which school team do you mean?

\vspace{0.5em}
\noindent \textbf{Speaker 1 [00:22]:} The football school team.

\vspace{0.5em}
\noindent \textbf{Speaker 2 [00:22]:} Which football school team? Junior high or senior high?

\vspace{0.5em}
\noindent \textbf{Speaker 1 [00:25]:} Senior high. In your typical team, what position do you usually play? Mainly?

\vspace{0.5em}
\noindent \textbf{Speaker 2 [00:30]:} Full-back, generally right-back, I play right-back more often.

\vspace{0.5em}
\noindent \textbf{Speaker 1 [00:36]:} Okay, right-back. How do you define your responsibilities on the pitch? What do you usually do when you are in this position?

\vspace{0.5em}
\noindent \textbf{Speaker 2 [00:53]:} You have to separate it into offense and defense. For offense, when I receive the ball on this side, my first choice is to give it to the midfield. Okay. Then, if it's a fast break, doing a diagonal long pass to switch it to the other side, giving it to the winger on the other side. Also, either moving down the flank to give it to the winger on this side, or I move up myself. Okay. Defense depends on which side they come from. Defense also depends on the position, and offense depends on the position too. In the backfield during offense, it's mainly about circulating the ball. If it's in the midfield or frontfield, you can carry it if you want to, or try a couple of shots yourself, if offense is in the opponent's half...

\vspace{0.5em}
\noindent \textbf{Speaker 2 [01:49]:} Wait, defense, we are the defensive side, they are the opponent. When the defensive side is in the opponent's half, it mainly depends on whether teammates are in position. If teammates have dropped back into position, you can move up and challenge for the ball boldly. If teammates are not in position, it's the same for all positions: try to delay the attack to give teammates time to drop back into position. When the defense reaches the midfield, try to force the attacking players to the flanks. After forcing them to the flanks at midfield, you can use your body to dispossess them and intercept the ball.

\vspace{0.5em}
\noindent \textbf{Speaker 2 [02:37]:} Reaching the backfield, in the backfield, if he cuts inside, you have to stabilize your center of gravity because there are more tricks and movements in front of the penalty area. If he fakes outward and crosses from the baseline, you have to watch the timing of his cross. Also, playing as a full-back requires listening to the center-back's instructions, being careful not to lag behind when trapping opponents offside, failing the offside trap and letting them steal a goal. Communication in the backline is very important; you communicate who covers and who defends. If your position is exposed and the center-back goes to cover your position, then you have to run to the center to cover for the center-back. Throw-ins depend on the tactics. If it's a long pass from the goalkeeper looking for the front players, then tuck inside. If it's about circulating the ball, then pull as close to the touchline as possible. What you mainly do also depends on the tactics, for example, Total Football, or counter-attacking. If playing an open attacking game, when attacking you can move up and become a winger. The winger tucks inside and becomes two defensive midfielders. This is for 8-man football; 11-man football should be similar in principle. If it's a counter-attack, just stay in your position well, and when you have the ball, look for the players upfront.

\vspace{0.5em}
\noindent \textbf{Speaker 1 [04:13]:} Mhm, okay, okay. Do you still remember the last match you played? When you played as a backfield player, how much of a decisive role do you think you played in your team's victory or defeat?

\vspace{0.5em}
\noindent \textbf{Speaker 2 [04:28]:} Playing the backline in the last match? Yes.

\vspace{0.5em}
\noindent \textbf{Speaker 2 [04:38]:} No, then...

\vspace{0.5em}
\noindent \textbf{Speaker 1 [04:39]:} Do you think your role was large? Compared to the frontfield players?

\vspace{0.5em}
\noindent \textbf{Speaker 2 [04:47]:} There's nothing to compare. The frontfield looks for chances to score, and the backfield wins the ball back to create attacking opportunities for the frontfield. I think this impact, unless it's a goal-line clearance or a solo run, which could be said to have a larger impact, in a normal match, I think the impact is about the same.

\vspace{0.5em}
\noindent \textbf{Speaker 1 [05:11]:} Okay, when your mistake, for instance, leads to your team conceding a goal, do you feel a lot of pressure from teammates, coaches, or spectators? For example, a direct mistake.

\vspace{0.5em}
\noindent \textbf{Speaker 2 [05:23]:} That depends on which team it is, and also on what... generally like primary and junior high school teams, they would...

\vspace{0.5em}
\noindent \textbf{Speaker 2 [05:37]:} Generally, my dad and the coach would pressure me, saying some harsh words. Later, I went to another club team where people were exceptionally nice, and on the pitch, we players would also encourage each other. When conceding a goal, try not to complain, but focus on encouragement.

\vspace{0.5em}
\noindent \textbf{Speaker 1 [06:04]:} Then when the team wins, to what extent do you think this credit is attributed to the defense of the backfield players? And to what extent is it viewed as the result of the entire team's joint efforts?

\vspace{0.5em}
\noindent \textbf{Speaker 2 [06:14]:} I think in most cases it's the credit of the entire team's joint efforts. Individual extreme cases are not ruled out, for example, if the frontfield is exceptionally weak, or if a team has an exceptionally strong frontfield but an exceptionally weak backfield. In such extremely extreme situations, one has to judge who made a greater contribution based on the circumstances at the time.

\vspace{0.5em}
\noindent \textbf{Speaker 1 [06:38]:} Do you think defensive players face harsher criticism or negative attention than offensive players? For instance, a frontfield player made a mistake but later scored a goal to level the score. Do you think this kind of situation...?

\vspace{0.5em}
\noindent \textbf{Speaker 2 [06:58]:} Then I think it's normal, because playing as a defender has requirements for the mistake rate. But for a forward, if they make a slight mistake, they can still chase back. If you make a mistake in the backfield as a defender, it's very likely to result in a goal. So playing as a defender, there is negative attention, but the main thing is to stay steady yourself. You can't let yourself get flustered because of sarcastic remarks from the sidelines, which would lead to worse performance later and more missed balls, creating a vicious cycle. Sometimes negative evaluations come from the opponents, and in those cases, you can completely ignore them. Frontfield players... because my team basically played counter-attacks, our frontfield players relied on personal ability. But the frontfield and midfield shouldn't just make random mistakes either.

\vspace{0.5em}
\noindent \textbf{Speaker 1 [08:01]:} Okay, then for example, during post-match summaries, or when you usually discuss, do you more frequently discuss goals or assists, or a defender's passing under pressure, or that kind of critical sliding tackle?

\vspace{0.5em}
\noindent \textbf{Speaker 2 [08:17]:} We don't discuss those; it would mostly be normal post-match reviews, which review the rhythm of the entire match.

\vspace{0.5em}
\noindent \textbf{Speaker 1 [08:27]:} In daily discussions among classmates, do goals and assists become more of a topic, or do sliding tackles and clearances tend to become the topic?

\vspace{0.5em}
\noindent \textbf{Speaker 2 [08:39]:} Become a topic.

\vspace{0.5em}
\noindent \textbf{Speaker 1 [08:41]:} Meaning more...

\vspace{0.5em}
\noindent \textbf{Speaker 2 [08:43]:} Then I think it depends on the situation. For example, if the team is trailing and you score a goal in a very beautiful way, I think that will become a topic. Also, if it's a draw and a defender makes a goal-line clearance, that will also become a topic. I think it's actually about the same---the highlights of the match.

\vspace{0.5em}
\noindent \textbf{Speaker 1 [09:14]:} Okay, then do you feel your defensive contribution, like your number of tackles or clearances, because it's hard to be directly quantified like goals... For instance, we often know how many goals a match had, but few people count or can easily count your clearance frequency or sliding tackle frequency, making it hard to directly quantify. Because of this, do you feel your defensive contribution has been underestimated?

\vspace{0.5em}
\noindent \textbf{Speaker 2 [09:38]:} No, I think evaluating a player's contribution solely by quantification is not very reasonable. Because that would be equivalent to using dribbling time, passing accuracy, and take-on success rate to judge the contribution of a forward or winger, which I think is unreasonable. The evaluation of defensive players, I think, should be from an overall perspective. Were you steady in this match? Did you make many mistakes? I think your contribution rate can be judged from here. Did the defensive players create many opportunities for the frontfield players? Was the connection---passing the ball to midfield players and connecting forward---efficient? I think defensive players should be evaluated based on such an overall view, rather than through quantified data like sliding tackles and clearances.

\vspace{0.5em}
\noindent \textbf{Speaker 1 [10:43]:} You mentioned this overall evaluation, could you be more specific, or is there any method that can make this... because I feel there might be a risk of being undervalued. Do you think there are any measures, policies, or new methods to make this defensive contribution more visible? Can you think of a specific method or measure? It could be in the school team, or for example, updates to data statistics in international tournaments, or regular education in schools, or policies on online platforms.

\vspace{0.5em}
\noindent \textbf{Speaker 2 [11:43]:} Players on the pitch... it's actually normal for players' contributions on the pitch to be undervalued.

\vspace{0.5em}
\noindent \textbf{Speaker 1 [11:50]:} Very normal.

\vspace{0.5em}
\noindent \textbf{Speaker 2 [11:51]:} Yeah, it can't be solved right now, because there's a portion of the football-watching crowd who don't play football. It's not that they don't understand football, but they haven't been on the pitch, so they don't understand that feeling, and therefore they will focus more on goals---the highlight performances of the offensive side. So the defensive side will definitely be undervalued. But there will still be people who think you play well; gold will always shine. If you play well enough on defense, people will still notice your contribution. But the premise is that you make certain contributions on both ends of the pitch for it to be relatively obvious. If you are purely defensive, that probably won't work well. If you are purely defensive or purely offensive, if you don't defend, people will scold you too. If you are involved in both offense and defense, you should be noticed.

\vspace{0.5em}
\noindent \textbf{Speaker 1 [12:51]:} So you mean backfield players can participate more in the offense without losing the overall strength of their defense. Yes.

\vspace{0.5em}
\noindent \textbf{Speaker 2 [13:00]:} And then, the two full-backs of the defensive line can frequently move forward to join the attack. Then... how should I say it? It depends on the team's overall strength and your personal ability. If you can go up and make it back, then moving up is definitely good. If you can't make it back, given limited ability, my advice is to try your best to defend. At this time, if your contribution is undervalued, then it shouldn't be called undervalued; it's just that your contribution is indeed relatively small.

\vspace{0.5em}
\noindent \textbf{Speaker 1 [13:32]:} Oh, okay.

\vspace{0.5em}
\noindent \textbf{Speaker 2 [13:33]:} Offensive players also sometimes come back to participate in defense, it's just relatively rare, and there's no requirement for them to absolutely come back and defend. But it's equivalent to saying that defense participating in offense is like adding a little bit on top of your duty, while offense coming back to defend is like a bonus.

\vspace{0.5em}
\noindent \textbf{Speaker 1 [13:59]:} Then do you think this potential oversight---contribution being ignored---will affect status, self-confidence, or market value? Professional players.

\vspace{0.5em}
\noindent \textbf{Speaker 2 [14:10]:} I don't know about market value. Status definitely. Status definitely. It's like, definitely no one blames the goalkeeper.

\vspace{0.5em}
\noindent \textbf{Speaker 1 [14:16]:} Oh, right, yes. In the eyes of the general public?

\vspace{0.5em}
\noindent \textbf{Speaker 2 [14:26]:} In the eyes of the general public regarding what? Like status, strength, etc.

\vspace{0.5em}
\noindent \textbf{Speaker 1 [14:30]:} Because a defender's contribution is underestimated or ignored, do you think their potential status in the eyes of fans might decrease?

\vspace{0.5em}
\noindent \textbf{Speaker 2 [14:57]:} I don't know about that, because I'm not a professional player.

\vspace{0.5em}
\noindent \textbf{Speaker 1 [15:00]:} What about amateurs? Like in our school.

\vspace{0.5em}
\noindent \textbf{Speaker 2 [15:03]:} Relative to amateurs, I think it's definitely not undervalued. I think this phenomenon shouldn't exist.

\vspace{0.5em}
\noindent \textbf{Speaker 1 [15:10]:} Okay, then this concludes it. OK. Great, thank you.

\end{linenumbers}

\end{document}